\documentclass{ximera}

\input{../../../preamble.tex}


\definecolor{fvdnavy}{HTML}{071D43}
\definecolor{fvdcyan}{HTML}{20BFC6}
\definecolor{fvdcyanlight}{HTML}{EFFCFB}
\definecolor{fvdmagenta}{HTML}{F10D69}
\definecolor{fvdmagentalight}{HTML}{FFF1F7}
\definecolor{fvdtext}{HTML}{163044}





%=================================================
% Pedagogische omgevingen
% PDF: tcolorbox
% HTML: learning-box
%=================================================

\newenvironment{voorkennis}
{%
    \ifdefined\HCode
        \HCode{%
<section class="learning-box learning-box--prior">
<div class="learning-box__header">
<span class="learning-box__icon" aria-hidden="true"></span>
<span class="learning-box__title">Voorkennis</span>
</div>
<div class="learning-box__content">
        }%
        \begin{itemize}
    \else
       \begin{tcolorbox}[
    enhanced,
    breakable,
    colback=fvdcyanlight,
    colframe=fvdcyan,
    coltext=fvdtext,
    boxrule=0.5pt,
    leftrule=4pt,
    arc=4mm,
    outer arc=4mm,
    left=7mm,
    right=6mm,
    top=5mm,
    bottom=5mm,
    before skip=6mm,
    after skip=6mm,
    title=Voorkennis,
    coltitle=fvdcyan,
    fonttitle=\bfseries\Large,
    attach title to upper={\par\smallskip},
    boxed title style={
        empty,
        left=0mm,
        right=0mm,
        top=0mm,
        bottom=0mm
    },
    overlay={
        \draw[
            fvdcyan,
            line width=0.5pt
        ]
        ([xshift=42mm,yshift=-13mm]frame.north west)
        --
        ([xshift=-7mm,yshift=-13mm]frame.north east);
    }
]
        \begin{itemize}
    \fi
}
{%
    \end{itemize}
    \ifdefined\HCode
        \HCode{%
</div>
</section>
        }%
    \else
        \end{tcolorbox}
    \fi
}


\newenvironment{leerdoelen}
{%
    \ifdefined\HCode
        \HCode{%
<section class="learning-box learning-box--goals">
<div class="learning-box__header">
<span class="learning-box__icon" aria-hidden="true"></span>
<span class="learning-box__title">Leerdoelen</span>
</div>
<div class="learning-box__content">
        }%
        \begin{itemize}
    \else
\begin{tcolorbox}[
    enhanced,
    breakable,
    colback=fvdmagentalight,
    colframe=fvdmagenta,
    coltext=fvdtext,
    boxrule=0.5pt,
    leftrule=4pt,
    arc=4mm,
    outer arc=4mm,
    left=7mm,
    right=6mm,
    top=5mm,
    bottom=5mm,
    before skip=6mm,
    after skip=6mm,
    title=Leerdoelen,
    coltitle=fvdmagenta,
    fonttitle=\bfseries\Large,
    attach title to upper={\par\smallskip},
    boxed title style={
        empty,
        left=0mm,
        right=0mm,
        top=0mm,
        bottom=0mm
    },
    overlay={
        \draw[
            fvdmagenta,
            line width=0.5pt
        ]
        ([xshift=42mm,yshift=-13mm]frame.north west)
        --
        ([xshift=-7mm,yshift=-13mm]frame.north east);
    }
]
        \begin{itemize}
    \fi
}
{%
    \end{itemize}
    \ifdefined\HCode
        \HCode{%
</div>
</section>
        }%
    \else
        \end{tcolorbox}
    \fi
}


\addPrintStyle{../../..}

\title{Oefeningen — Verandering en het afgeleid getal}
\author{Ingmar Herreman}


% ============================================================
% PDF: A4 LIGGEND
% ============================================================


\begin{document}

% FVD_CHAPTER_DATA_START
% Automatisch gegenereerd uit index.tex.
% Niet handmatig aanpassen.
\fvdchapterdata{1}{Veranderingen in beeld}{1}
% FVD_CHAPTER_DATA_END
























































































\maketitle

\FVDbeginExercises


% ============================================================
% BASIS
% ============================================================


% ------------------------------------------------------------
% OEFENING 1
% ------------------------------------------------------------

\begin{oefening}[Gemiddelde verandering]{\oefB\ESS}

Gegeven is

\[
f(x)=x^2-2x+3.
\]

\begin{question}
Bereken $f(1)$ en $f(4)$.
\end{question}

\begin{question}
Bereken de gemiddelde verandering van $f$ op $[1,4]$:

\[
\frac{f(4)-f(1)}{4-1}.
\]
\end{question}

\begin{question}
Welke meetkundige betekenis heeft dit resultaat?

\begin{multipleChoice}
\choice{De richtingscoëfficiënt van de raaklijn in $x=1$.}
\choice[correct]{De richtingscoëfficiënt van de secant door de punten met abscissen $1$ en $4$.}
\choice{De functiewaarde in $x=4$.}
\choice{De oppervlakte onder de grafiek tussen $1$ en $4$.}
\end{multipleChoice}
\end{question}

\begin{oplossing}

We berekenen eerst de functiewaarden:
\[
f(1)=1^2-2\cdot1+3=2,
\]
en
\[
f(4)=4^2-2\cdot4+3=11.
\]

De gemiddelde verandering op $[1,4]$ is dus
\[
\frac{f(4)-f(1)}{4-1}
=
\frac{11-2}{3}
=
3.
\]

Meetkundig is dit de richtingscoëfficiënt van de secant door de punten
\[
(1,2)
\qquad\text{en}\qquad
(4,11).
\]

\end{oplossing}

\end{oefening}


% ------------------------------------------------------------
% OEFENING 2
% ------------------------------------------------------------

\begin{oefening}[Van tabel naar verandering]{\oefB\ESS}

De positie $s$ van een bewegend voorwerp wordt op enkele tijdstippen
gemeten.

\[
\begin{array}{c|ccccc}
t\;(\mathrm{s}) & 0 & 1 & 2 & 3 & 4\\ \hline
s\;(\mathrm{m}) & 2 & 5 & 10 & 17 & 26
\end{array}
\]

\begin{question}
Bereken de gemiddelde snelheid tussen $t=1$ en $t=3$.
\end{question}

\begin{question}
Noteer de correcte eenheid.
\end{question}

\begin{question}
Bereken ook de gemiddelde snelheid tussen $t=2$ en $t=4$.
\end{question}

\begin{question}
Leg uit waarom een gemiddelde snelheid niet noodzakelijk gelijk is
aan de snelheid op elk tijdstip van het onderzochte interval.
\end{question}

\begin{oplossing}

Tussen $t=1$ en $t=3$ geldt
\[
v_{\mathrm{gem}}
=
\frac{s(3)-s(1)}{3-1}
=
\frac{17-5}{2}
=
6\ \mathrm{m/s}.
\]

De correcte eenheid is dus $\mathrm{m/s}$.

Tussen $t=2$ en $t=4$ krijgen we
\[
v_{\mathrm{gem}}
=
\frac{s(4)-s(2)}{4-2}
=
\frac{26-10}{2}
=
8\ \mathrm{m/s}.
\]

Een gemiddelde snelheid beschrijft de totale positieverandering gedeeld
door de totale tijdsduur van een interval. Ze zegt niet dat het voorwerp
op elk afzonderlijk tijdstip van dat interval precies die snelheid heeft.

\end{oplossing}

\end{oefening}


% ------------------------------------------------------------
% OEFENING 3
% ------------------------------------------------------------

\begin{oefening}[Het differentiequotiënt]{\oefB\ESS}

Gegeven is

\[
f(x)=x^2+1.
\]

We onderzoeken de verandering in de buurt van $x=2$.

\begin{question}
Bereken
\[
f(2).
\]
\end{question}

\begin{question}
Bereken en vereenvoudig
\[
\frac{f(x)-f(2)}{x-2}.
\]
\end{question}

\begin{question}
Bepaal vervolgens
\[
\lim_{x\to2}
\frac{f(x)-f(2)}{x-2}.
\]
\end{question}

\begin{question}
Welke betekenis heeft het gevonden getal?
\end{question}

\begin{oplossing}

We hebben
\[
f(2)=2^2+1=5.
\]

Volgens de definitie:
\[
f'(2)
=
\lim_{x\to2}
\frac{f(x)-f(2)}{x-2}.
\]

We vereenvoudigen eerst het differentiequotiënt:
\[
\begin{aligned}
\frac{f(x)-f(2)}{x-2}
&=
\frac{x^2+1-5}{x-2}\\
&=
\frac{x^2-4}{x-2}\\
&=
\frac{(x-2)(x+2)}{x-2}\\
&=
x+2
\qquad (x\neq2).
\end{aligned}
\]

Daarom
\[
f'(2)
=
\lim_{x\to2}(x+2)
=
\boxed{4}.
\]

Het getal $4$ is de ogenblikkelijke verandering van $f$ bij $x=2$.
Meetkundig is het de richtingscoëfficiënt van de raaklijn aan de grafiek
in het punt met abscis $2$.

\end{oplossing}

\end{oefening}


% ------------------------------------------------------------
% OEFENING 4
% ------------------------------------------------------------

\begin{oefening}[Afgeleid getal vanuit de definitie]{\oefB\ESS}

Gegeven is

\[
f(x)=3x^2-2x.
\]

Bereken vanuit de definitie

\[
f'(1)
=
\lim_{x\to1}
\frac{f(x)-f(1)}{x-1}.
\]

Werk het differentiequotiënt volledig uit.

\begin{oplossing}

Eerst:
\[
f(1)=3\cdot1^2-2\cdot1=1.
\]

We gebruiken de definitie
\[
f'(1)
=
\lim_{x\to1}
\frac{f(x)-f(1)}{x-1}.
\]

Het differentiequotiënt wordt
\[
\begin{aligned}
\frac{f(x)-f(1)}{x-1}
&=
\frac{3x^2-2x-1}{x-1}\\
&=
\frac{(x-1)(3x+1)}{x-1}\\
&=
3x+1
\qquad (x\neq1).
\end{aligned}
\]

Dus
\[
f'(1)
=
\lim_{x\to1}(3x+1)
=
\boxed{4}.
\]

\end{oplossing}

\end{oefening}


% ------------------------------------------------------------
% OEFENING 5
% ------------------------------------------------------------

\begin{oefening}[Wat betekent $f'(a)$?]{\oefB\ESS}

Duid telkens alle correcte interpretaties van $f'(a)$ aan.

\begin{selectAll}
\choice[correct]{De ogenblikkelijke verandering van $f$ in $a$.}
\choice[correct]{De richtingscoëfficiënt van de raaklijn aan de grafiek in $x=a$.}
\choice{De functiewaarde $f(a)$.}
\choice{De gemiddelde verandering van $f$ op elk interval dat $a$ bevat.}
\choice[correct]{De limiet van een differentiequotiënt, wanneer die limiet bestaat.}
\end{selectAll}

\begin{oplossing}

De correcte interpretaties zijn:

\begin{itemize}
\item $f'(a)$ is de ogenblikkelijke verandering van $f$ in $a$;
\item $f'(a)$ is de richtingscoëfficiënt van de raaklijn aan de grafiek
      in $x=a$;
\item $f'(a)$ is, wanneer de limiet bestaat, de limiet van een
      differentiequotiënt.
\end{itemize}

$f'(a)$ is niet hetzelfde als de functiewaarde $f(a)$ en ook niet als de
gemiddelde verandering op een willekeurig interval rond $a$.

\end{oplossing}

\end{oefening}


% ------------------------------------------------------------
% OEFENING 6
% ------------------------------------------------------------

\begin{oefening}[Teken van het afgeleid getal]{\oefB\ESS}

Beantwoord zonder een afgeleide te berekenen.

\begin{question}
De grafiek van $f$ heeft in $x=2$ een stijgende raaklijn.

Welk teken heeft $f'(2)$?
\end{question}

\begin{question}
De grafiek van $g$ heeft in $x=-1$ een dalende raaklijn.

Welk teken heeft $g'(-1)$?
\end{question}

\begin{question}
De grafiek van $h$ heeft in $x=4$ een horizontale raaklijn.

Wat weet je over $h'(4)$?
\end{question}

\begin{oplossing}

Een stijgende raaklijn heeft een positieve richtingscoëfficiënt. Dus
\[
f'(2)>0.
\]

Een dalende raaklijn heeft een negatieve richtingscoëfficiënt. Dus
\[
g'(-1)<0.
\]

Een horizontale raaklijn heeft richtingscoëfficiënt nul. Dus
\[
h'(4)=0.
\]

\end{oplossing}

\end{oefening}


% ------------------------------------------------------------
% OEFENING 7
% ------------------------------------------------------------

\begin{oefening}[Raaklijn bepalen]{\oefB\ESS}

Gegeven is

\[
f(x)=x^2-3x+4
\]

en

\[
f'(2)=1.
\]

\begin{question}
Bepaal het raakpunt aan de grafiek voor $x=2$.
\end{question}

\begin{question}
Welke richtingscoëfficiënt heeft de raaklijn?
\end{question}

\begin{question}
Bepaal een vergelijking van de raaklijn.
\end{question}

\begin{oplossing}

Het raakpunt heeft abscis $2$. De bijbehorende functiewaarde is
\[
f(2)=2^2-3\cdot2+4=2.
\]

Het raakpunt is dus
\[
P(2,2).
\]

Omdat
\[
f'(2)=1,
\]
heeft de raaklijn richtingscoëfficiënt $1$.

Met de punt-richtingsvorm:
\[
y-2=1(x-2).
\]

Dus
\[
\boxed{y=x}.
\]

\end{oplossing}

\end{oefening}


% ------------------------------------------------------------
% OEFENING 8
% ------------------------------------------------------------

\begin{oefening}[Van raaklijn naar afgeleid getal]{\oefB}

De raaklijn aan de grafiek van $f$ in het punt $P(3,5)$ heeft
vergelijking

\[
y=-2x+11.
\]

\begin{question}
Bepaal $f(3)$.
\end{question}

\begin{question}
Bepaal $f'(3)$.
\end{question}

\begin{question}
Is $f$ in de onmiddellijke omgeving van $x=3$ stijgend of dalend?
Motiveer.
\end{question}

\begin{oplossing}

Omdat $P(3,5)$ op de grafiek van $f$ ligt,
\[
\boxed{f(3)=5}.
\]

De raaklijn heeft vergelijking
\[
y=-2x+11,
\]
en dus richtingscoëfficiënt $-2$. Bijgevolg
\[
\boxed{f'(3)=-2}.
\]

De raaklijn is dalend. De grafiek heeft bij $x=3$ dus een negatieve
ogenblikkelijke verandering.

\end{oplossing}

\end{oefening}


% ------------------------------------------------------------
% OEFENING 9
% ------------------------------------------------------------

\begin{oefening}[Eenheden]{\oefB\ESS}

De hoogte $h(t)$ van een object wordt uitgedrukt in meter en de tijd
$t$ in seconde.

\begin{question}
Welke eenheid heeft

\[
\frac{h(5)-h(2)}{5-2}?
\]
\end{question}

\begin{question}
Welke eenheid heeft

\[
h'(5)?
\]
\end{question}

\begin{question}
Leg uit waarom beide grootheden dezelfde eenheid hebben, hoewel hun
betekenis verschillend is.
\end{question}

\begin{oplossing}

De teller $h(5)-h(2)$ wordt uitgedrukt in meter en de noemer $5-2$
in seconde. Daarom heeft
\[
\frac{h(5)-h(2)}{5-2}
\]
als eenheid
\[
\boxed{\mathrm{m/s}}.
\]

Ook
\[
h'(5)
\]
heeft als eenheid
\[
\boxed{\mathrm{m/s}}.
\]

Beide grootheden beschrijven een verandering van hoogte per tijdseenheid.
Het verschil zit in de betekenis: het differentiequotiënt geeft een
gemiddelde verandering over een interval, terwijl $h'(5)$ de
ogenblikkelijke verandering bij $t=5$ beschrijft.

\end{oplossing}

\end{oefening}


% ------------------------------------------------------------
% OEFENING 10
% ------------------------------------------------------------

\begin{oefening}[Horizontale raaklijn]{\oefB\ESS}

Een leerling beweert:

\[
f'(a)=0
\quad\Longrightarrow\quad
f \text{ heeft in }a\text{ een extremum}.
\]

\begin{question}
Is deze uitspraak altijd juist?
\end{question}

\begin{question}
Geef een tegenvoorbeeld of beschrijf een mogelijke grafiek waarvoor
$f'(a)=0$ maar $f$ geen extremum heeft in $a$.
\end{question}

\begin{question}
Wat kun je uit $f'(a)=0$ wél besluiten?
\end{question}

\begin{oplossing}

De uitspraak is niet altijd juist.

Een eenvoudig tegenvoorbeeld is
\[
f(x)=x^3
\]
in $a=0$. Daar geldt
\[
f'(0)=0,
\]
maar de functie blijft zowel links als rechts van $0$ stijgen. Er is dus
geen maximum of minimum.

Uit
\[
f'(a)=0
\]
kunnen we wel besluiten dat de raaklijn in $x=a$ horizontaal is.
Om een extremum te besluiten, moeten we bovendien het verloop of het
teken van de afgeleide links en rechts van $a$ onderzoeken.

\end{oplossing}

\end{oefening}


% ============================================================
% VERDIEPING
% ============================================================




% ------------------------------------------------------------
% OEFENING 11
% ------------------------------------------------------------

\begin{oefening}[Secanten naderen een raaklijn]{\oefU\ESS}

Gegeven is

\[
f(x)=x^2.
\]

We onderzoeken het punt met abscis $x=1$.

Vul de tabel aan.

\[
\begin{array}{c|c}
x &
\dfrac{f(x)-f(1)}{x-1}
\\ \hline
2 & \\
1.5 & \\
1.1 & \\
1.01 & \\
0.9 & \\
0.99 &
\end{array}
\]

\begin{question}
Naar welk getal lijken de differentiequotiënten te naderen?
\end{question}

\begin{question}
Wat betekent dit getal meetkundig?
\end{question}

\begin{question}
Waarom onderzoeken we waarden van $x$ zowel links als rechts van $1$?
\end{question}

\begin{oplossing}

Voor
\[
f(x)=x^2
\]
geldt
\[
\frac{f(x)-f(1)}{x-1}
=
\frac{x^2-1}{x-1}
=
x+1
\qquad (x\neq1).
\]

Daarmee krijgen we:
\[
\begin{array}{c|c}
x & \dfrac{f(x)-f(1)}{x-1}\\ \hline
2 & 3\\
1.5 & 2.5\\
1.1 & 2.1\\
1.01 & 2.01\\
0.9 & 1.9\\
0.99 & 1.99
\end{array}
\]

De waarden naderen naar
\[
\boxed{2}.
\]

Meetkundig is dit de richtingscoëfficiënt van de raaklijn aan de grafiek
in $x=1$.

We onderzoeken waarden van $x$ zowel links als rechts van $1$ omdat we
willen nagaan of de differentiequotiënten van beide kanten naar hetzelfde
getal naderen.

\end{oplossing}

\end{oefening}


% ------------------------------------------------------------
% OEFENING 12
% ------------------------------------------------------------

\begin{oefening}[Van grafiek naar afgeleide informatie]{\oefU\ESS}

Onderstaande figuur toont een mogelijke grafiek van $f$.

\begin{center}
\begin{tikzpicture}[x=1cm,y=0.75cm]

  \draw[->,draw=black!45] (-4.2,0) -- (4.5,0)
    node[right] {$x$};
  \draw[->,draw=black!45] (0,-3.2) -- (0,4.3)
    node[above] {$y$};

  \draw[
    draw=fvdnavy,
    very thick,
    smooth,
    domain=-3.4:3.4,
    samples=120
  ]
  plot (\x,{0.12*(\x)^3-1.1*(\x)});

  \fill[fill=fvdmagenta] (-2,1.36) circle (2.2pt);
  \fill[fill=fvdmagenta] (0,0) circle (2.2pt);
  \fill[fill=fvdmagenta] (2,-1.36) circle (2.2pt);

  \node[below] at (-2,0) {$a$};
  \node[below] at (0,0) {$b$};
  \node[above] at (2,0) {$c$};

\end{tikzpicture}
\end{center}

\begin{question}
Bepaal op basis van de grafiek het teken van

\[
f'(a),\qquad f'(b),\qquad f'(c).
\]

Je hoeft geen exacte waarden te bepalen.
\end{question}

\begin{question}
In welk van de drie punten lijkt de raaklijn het steilst?
\end{question}

\begin{question}
Wat zegt een grotere absolute waarde van $f'(x)$ over de raaklijn?
\end{question}

\begin{oplossing}

In het punt met abscis $a$ stijgt de grafiek. Daarom
\[
f'(a)>0.
\]

In het punt met abscis $b$ daalt de grafiek. Daarom
\[
f'(b)<0.
\]

In het punt met abscis $c$ stijgt de grafiek opnieuw. Daarom
\[
f'(c)>0.
\]

Op basis van de figuur lijkt de raaklijn in $b$ het steilst.

Hoe groter de absolute waarde
\[
|f'(x)|,
\]
hoe steiler de raaklijn in dat punt. Het teken van $f'(x)$ bepaalt of die
raaklijn stijgt of daalt.

\end{oplossing}

\end{oefening}


% ------------------------------------------------------------
% OEFENING 13
% ------------------------------------------------------------

\begin{oefening}[Raaklijn uit de definitie]{\oefU\ESS}

Gegeven is

\[
f(x)=2x^2-x+1.
\]

Bepaal de vergelijking van de raaklijn aan de grafiek in het punt
met abscis $x=2$.

Je mag geen afleidingsregel gebruiken.

Je oplossing moet dus de volgende keten bevatten:

\[
\frac{f(x)-f(2)}{x-2}
\longrightarrow
f'(2)
\longrightarrow
\text{raaklijn}.
\]

\begin{oplossing}

Eerst bepalen we
\[
f(2)=2\cdot2^2-2+1=7.
\]

Volgens de definitie:
\[
f'(2)
=
\lim_{x\to2}
\frac{f(x)-f(2)}{x-2}.
\]

We vereenvoudigen:
\[
\begin{aligned}
\frac{f(x)-f(2)}{x-2}
&=
\frac{2x^2-x+1-7}{x-2}\\
&=
\frac{2x^2-x-6}{x-2}\\
&=
\frac{(x-2)(2x+3)}{x-2}\\
&=
2x+3
\qquad (x\neq2).
\end{aligned}
\]

Daaruit volgt
\[
f'(2)
=
\lim_{x\to2}(2x+3)
=
7.
\]

De raaklijn gaat door $(2,7)$ en heeft richtingscoëfficiënt $7$:
\[
y-7=7(x-2).
\]

Dus
\[
\boxed{y=7x-7}.
\]

\end{oplossing}

\end{oefening}


% ------------------------------------------------------------
% OEFENING 14
% ------------------------------------------------------------

\begin{oefening}[Snelheid op één ogenblik]{\oefU\ESS}

De positie van een bewegend voorwerp wordt beschreven door

\[
s(t)=t^2+2t
\]

met $s$ in meter en $t$ in seconde.

\begin{question}
Bereken de gemiddelde snelheid tussen $t=2$ en een tijdstip $t\neq2$:
\[
\frac{s(t)-s(2)}{t-2}.
\]
\end{question}

\begin{question}
Vereenvoudig je uitdrukking.
\end{question}

\begin{question}
Laat $t\to2$ en bepaal de ogenblikkelijke snelheid op $t=2$.
\end{question}

\begin{question}
Interpreteer je resultaat in een volledige zin met eenheid.
\end{question}

\begin{oplossing}

We berekenen eerst
\[
s(2)=2^2+2\cdot2=8.
\]

De gemiddelde snelheid tussen $t=2$ en een tijdstip $t\neq2$ is
\[
\begin{aligned}
\frac{s(t)-s(2)}{t-2}
&=
\frac{t^2+2t-8}{t-2}\\
&=
\frac{(t-2)(t+4)}{t-2}\\
&=
t+4.
\end{aligned}
\]

Voor $t\to2$ krijgen we
\[
s'(2)
=
\lim_{t\to2}
\frac{s(t)-s(2)}{t-2}
=
\lim_{t\to2}(t+4)
=
\boxed{6}.
\]

De ogenblikkelijke snelheid op $t=2$ s bedraagt dus
\[
\boxed{6\ \mathrm{m/s}}.
\]

Op het tijdstip $t=2$ s verandert de positie ogenblikkelijk met
$6$ meter per seconde in de positieve richting.

\end{oplossing}

\end{oefening}


% ------------------------------------------------------------
% OEFENING 15
% ------------------------------------------------------------

\begin{oefening}[Temperatuur verandert]{\oefU}

De temperatuur van een vloeistof wordt gedurende een experiment
beschreven door een functie $T(t)$, met $T$ in $^\circ\mathrm{C}$
en $t$ in minuten.

Op een bepaald ogenblik geldt

\[
T'(6)=-1.8.
\]

\begin{question}
Wat betekent het minteken?
\end{question}

\begin{question}
Wat betekent de waarde $1.8$?
\end{question}

\begin{question}
Geef een correcte interpretatie van

\[
T'(6)=-1.8
\]

in de context.
\end{question}

\begin{question}
Kun je hieruit besluiten dat de temperatuur gedurende de volledige
zesde minuut exact $1.8^\circ\mathrm{C}$ daalt? Verklaar.
\end{question}

\begin{oplossing}

Het minteken in
\[
T'(6)=-1.8
\]
betekent dat de temperatuur op dat ogenblik daalt.

De waarde $1.8$ geeft aan hoe snel de temperatuur op dat ogenblik
verandert: de grootte van de ogenblikkelijke verandering is
$1.8^\circ\mathrm{C}$ per minuut.

Een volledige interpretatie is:

\[
\boxed{\text{Na 6 minuten daalt de temperatuur ogenblikkelijk met }
1.8^\circ\mathrm{C/min}.}
\]

Hieruit volgt niet dat de temperatuur gedurende een volledige minuut
exact $1.8^\circ\mathrm{C}$ daalt. $T'(6)$ beschrijft alleen de
ogenblikkelijke verandering bij $t=6$.

\end{oplossing}

\end{oefening}


% ------------------------------------------------------------
% OEFENING 16
% ------------------------------------------------------------

\begin{oefening}[Welke informatie is voldoende?]{\oefU\ESS}

We willen $f'(3)$ bepalen.

Onderzoek telkens of de gegeven informatie voldoende is.

\begin{question}
Je kent alleen $f(3)=7$.
\end{question}

\begin{question}
Je kent twee punten van de grafiek: $(2,4)$ en $(4,8)$.
\end{question}

\begin{question}
Je kent de vergelijking van de raaklijn aan de grafiek in $x=3$.
\end{question}

\begin{question}
Je kent voor alle $x\neq3$ de uitdrukking

\[
\frac{f(x)-f(3)}{x-3}=2x-1.
\]

Bepaal in dit geval ook $f'(3)$.
\end{question}

\begin{oplossing}

\textbf{1. Alleen $f(3)=7$: niet voldoende.}

De functiewaarde vertelt waar het punt op de grafiek ligt, maar niet hoe
steil de grafiek daar is.

\medskip

\textbf{2. De punten $(2,4)$ en $(4,8)$: niet voldoende.}

Daarmee kunnen we wel de richtingscoëfficiënt van een secant bepalen:
\[
\frac{8-4}{4-2}=2.
\]
Dit is een gemiddelde verandering en hoeft niet gelijk te zijn aan
$f'(3)$.

\medskip

\textbf{3. De vergelijking van de raaklijn in $x=3$: voldoende.}

De richtingscoëfficiënt van die raaklijn is precies $f'(3)$.

\medskip

\textbf{4. Het differentiequotiënt is $2x-1$: voldoende.}

We nemen de limiet:
\[
f'(3)
=
\lim_{x\to3}
\frac{f(x)-f(3)}{x-3}
=
\lim_{x\to3}(2x-1)
=
\boxed{5}.
\]

\end{oplossing}

\end{oefening}


% ------------------------------------------------------------
% OEFENING 17
% ------------------------------------------------------------

\begin{oefening}[Een fout opsporen]{\oefU\ESS}

Een leerling wil voor

\[
f(x)=x^2+3
\]

het afgeleid getal in $x=2$ berekenen en schrijft:

\[
f'(2)
=
\lim_{x\to2}
\frac{f(x)}{x-2}
=
\lim_{x\to2}
\frac{x^2+3}{x-2}.
\]

\begin{question}
Waar zit de eerste fout?
\end{question}

\begin{question}
Schrijf het correcte differentiequotiënt.
\end{question}

\begin{question}
Bereken vervolgens $f'(2)$ correct.
\end{question}

\begin{question}
Leg uit waarom de term $-f(2)$ essentieel is in het
differentiequotiënt.
\end{question}

\begin{oplossing}

De eerste fout zit in het differentiequotiënt. In de teller moet de
verandering van de functiewaarde staan:
\[
f(x)-f(2).
\]

Omdat
\[
f(2)=2^2+3=7,
\]
is het correcte differentiequotiënt
\[
\frac{f(x)-f(2)}{x-2}
=
\frac{x^2+3-7}{x-2}.
\]

Uitwerken geeft
\[
\frac{x^2-4}{x-2}
=
\frac{(x-2)(x+2)}{x-2}
=
x+2
\qquad (x\neq2).
\]

Dus
\[
f'(2)
=
\lim_{x\to2}(x+2)
=
\boxed{4}.
\]

De term $-f(2)$ is essentieel omdat het differentiequotiënt de
\emph{verandering} van de functiewaarde meet. Daardoor nadert de teller
ook naar nul wanneer $x$ naar $2$ nadert.

\end{oplossing}

\end{oefening}


% ============================================================
% GEVORDERD
% ============================================================



% ------------------------------------------------------------
% OEFENING 18
% ------------------------------------------------------------

\begin{oefening}[Een onbekende parameter]{\oefG}

Gegeven is de familie functies

\[
f_a(x)=x^2+ax.
\]

De raaklijn aan de grafiek in het punt met abscis $x=1$ heeft
richtingscoëfficiënt $5$.

Bepaal $a$ vanuit de definitie van het afgeleid getal.

\begin{oplossing}

We gebruiken de definitie van het afgeleid getal:
\[
f_a'(1)
=
\lim_{x\to1}
\frac{f_a(x)-f_a(1)}{x-1}.
\]

Voor
\[
f_a(x)=x^2+ax
\]
geldt
\[
f_a(1)=1+a.
\]

Dus
\[
\begin{aligned}
\frac{f_a(x)-f_a(1)}{x-1}
&=
\frac{x^2+ax-(1+a)}{x-1}\\
&=
\frac{(x-1)(x+1+a)}{x-1}\\
&=
x+1+a
\qquad (x\neq1).
\end{aligned}
\]

Daarom
\[
f_a'(1)
=
\lim_{x\to1}(x+1+a)
=
2+a.
\]

De richtingscoëfficiënt is gegeven als $5$, dus
\[
2+a=5.
\]

Hieruit volgt
\[
\boxed{a=3}.
\]

\end{oplossing}

\end{oefening}


% ------------------------------------------------------------
% OEFENING 19
% ------------------------------------------------------------

\begin{oefening}[Zelfde afgeleid getal, andere functie]{\oefG}

Zoek twee verschillende veeltermfuncties $f$ en $g$ waarvoor

\[
f(2)\neq g(2)
\]

maar

\[
f'(2)=g'(2)=3.
\]

\begin{question}
Geef twee mogelijke voorschriften.
\end{question}

\begin{question}
Controleer de voorwaarden.
\end{question}

\begin{question}
Leg meetkundig uit hoe twee grafieken in verschillende punten toch
hetzelfde afgeleid getal kunnen hebben.
\end{question}

\begin{oplossing}

Een mogelijk voorbeeld is
\[
f(x)=3x
\qquad\text{en}\qquad
g(x)=3x+1.
\]

Dan
\[
f(2)=6
\qquad\text{en}\qquad
g(2)=7,
\]
zodat
\[
f(2)\neq g(2).
\]

Beide functies hebben overal richtingscoëfficiënt $3$, dus
\[
f'(2)=g'(2)=3.
\]

Meetkundig betekent dit dat de raaklijnen aan beide grafieken bij
$x=2$ dezelfde richtingscoëfficiënt hebben. De raaklijnen zijn dus
evenwijdig, hoewel de punten waarin ze de grafieken raken verschillend
kunnen liggen.

\end{oplossing}

\end{oefening}


% ------------------------------------------------------------
% OEFENING 20
% ------------------------------------------------------------

\begin{oefening}[Meetgegevens en een ogenblikkelijke snelheid]{\oefG}

Een sensor registreert de positie van een robot rond $t=4$ s.

\[
\begin{array}{c|ccccc}
t\;(\mathrm{s})
&3.8&3.9&4.0&4.1&4.2\\ \hline
s\;(\mathrm{m})
&7.22&7.60&8.00&8.42&8.86
\end{array}
\]

Het exacte functievoorschrift is onbekend.

\begin{question}
Bereken een benadering van de snelheid op $t=4$ door alleen de
metingen op $t=4$ en $t=4.1$ te gebruiken.
\end{question}

\begin{question}
Bereken een tweede benadering met de metingen op $t=3.9$ en $t=4$.
\end{question}

\begin{question}
Gebruik vervolgens de meetpunten op $t=3.9$ en $t=4.1$ om een
derde schatting te maken.
\end{question}

\begin{question}
Welke schatting vind je het meest overtuigend als benadering van de
ogenblikkelijke snelheid? Motiveer.
\end{question}

\begin{question}
Waarom kunnen we uit een eindige tabel met meetgegevens in het
algemeen niet de exacte waarde van $s'(4)$ bewijzen?
\end{question}

\begin{oplossing}

Met de metingen op $t=4$ en $t=4.1$:
\[
v
\approx
\frac{8.42-8.00}{4.1-4.0}
=
\frac{0.42}{0.1}
=
4.2\ \mathrm{m/s}.
\]

Met de metingen op $t=3.9$ en $t=4$:
\[
v
\approx
\frac{8.00-7.60}{4.0-3.9}
=
\frac{0.40}{0.1}
=
4.0\ \mathrm{m/s}.
\]

Met de punten symmetrisch rond $t=4$:
\[
v
\approx
\frac{8.42-7.60}{4.1-3.9}
=
\frac{0.82}{0.2}
=
4.1\ \mathrm{m/s}.
\]

De derde schatting,
\[
\boxed{4.1\ \mathrm{m/s}},
\]
is het meest overtuigend, omdat ze meetpunten gebruikt die even ver
links en rechts van $t=4$ liggen.

Uit een eindige tabel kunnen we in het algemeen de exacte waarde van
$s'(4)$ niet bewijzen. Voor een exacte afgeleide moeten we weten wat er
gebeurt wanneer het tijdsverschil willekeurig klein wordt; een eindige
tabel bevat slechts een beperkt aantal metingen.

\end{oplossing}

\end{oefening}


% ------------------------------------------------------------
% OEFENING 21
% ------------------------------------------------------------

\begin{oefening}[Een bewering onderzoeken]{\oefG}

Een leerling beweert:

\begin{quote}
Als twee functies in $x=a$ dezelfde functiewaarde hebben en
hetzelfde afgeleid getal, dan moeten hun grafieken in de buurt van
$a$ samenvallen.
\end{quote}

Onderzoek deze bewering.

\begin{question}
Is ze waar of onwaar?
\end{question}

\begin{question}
Geef een overtuigend voorbeeld of tegenvoorbeeld.
\end{question}

\begin{question}
Wat betekent het wél als

\[
f(a)=g(a)
\]

en

\[
f'(a)=g'(a)?
\]
\end{question}

\begin{oplossing}

De bewering is onwaar.

Neem bijvoorbeeld
\[
f(x)=x^2
\qquad\text{en}\qquad
g(x)=2x^2
\]
en kies $a=0$.

Dan
\[
f(0)=g(0)=0
\]
en
\[
f'(0)=g'(0)=0.
\]

Toch vallen de grafieken in de buurt van $0$ niet samen, want voor
$x\neq0$ geldt
\[
x^2\neq2x^2.
\]

Als
\[
f(a)=g(a)
\qquad\text{en}\qquad
f'(a)=g'(a),
\]
dan gaan beide grafieken door hetzelfde punt en hebben ze daar dezelfde
raaklijn. Dat betekent niet dat hun verdere verloop in de buurt van dat
punt hetzelfde moet zijn.

\end{oplossing}

\end{oefening}




\begin{oefening}[Alles samen]{\oefU\ESS}

Gegeven is

\[
f(x)=x^2-4x+5.
\]

Onderzoek het punt met abscis $x=3$.

\begin{question}
Bereken $f(3)$.
\end{question}

\begin{question}
Stel het differentiequotiënt

\[
\frac{f(x)-f(3)}{x-3}
\]

op en vereenvoudig.
\end{question}

\begin{question}
Bepaal vanuit de definitie $f'(3)$.
\end{question}

\begin{question}
Interpreteer het teken van $f'(3)$ grafisch.
\end{question}

\begin{question}
Bepaal de vergelijking van de raaklijn in $x=3$.
\end{question}

\begin{question}
Neem nu $x=3.1$.

Bereken de richtingscoëfficiënt van de overeenkomstige secant en
vergelijk die met $f'(3)$.
\end{question}

\begin{question}
Leg uit waarom die twee waarden dicht bij elkaar liggen.
\end{question}

\begin{oplossing}

Eerst:
\[
f(3)=3^2-4\cdot3+5=2.
\]

Volgens de definitie:
\[
f'(3)
=
\lim_{x\to3}
\frac{f(x)-f(3)}{x-3}.
\]

Het differentiequotiënt wordt
\[
\begin{aligned}
\frac{f(x)-f(3)}{x-3}
&=
\frac{x^2-4x+5-2}{x-3}\\
&=
\frac{x^2-4x+3}{x-3}\\
&=
\frac{(x-1)(x-3)}{x-3}\\
&=
x-1
\qquad (x\neq3).
\end{aligned}
\]

Daarom
\[
f'(3)
=
\lim_{x\to3}(x-1)
=
\boxed{2}.
\]

Omdat $f'(3)>0$, heeft de grafiek in $x=3$ een stijgende raaklijn.

Het raakpunt is $(3,2)$ en de richtingscoëfficiënt is $2$. Bijgevolg:
\[
y-2=2(x-3),
\]
dus
\[
\boxed{y=2x-4}.
\]

Voor $x=3.1$ is de richtingscoëfficiënt van de secant
\[
\frac{f(3.1)-f(3)}{3.1-3}
=
3.1-1
=
2.1.
\]

Deze waarde ligt dicht bij
\[
f'(3)=2,
\]
omdat het afgeleid getal de limiet is van de richtingscoëfficiënten van
de secanten wanneer $x$ naar $3$ nadert.

\end{oplossing}

\end{oefening}


\FVDendExercises
\end{document}